% ============================================================
% Chapter 2: Operating System Security: POSIX and Windows Architecture
% Language: Greek — edit freely; regenerate from src/content if preferred.
% ============================================================
\chapter{Ασφάλεια Λειτουργικών Συστημάτων: Αρχιτεκτονική POSIX και Windows}\label{ch:2}
\chaptermeta{Διακριτικός έλεγχος πρόσβασης · Ειδικά δικαιώματα · ACL και capabilities · Υποχρεωτικοί έλεγχοι · Token, UAC και επίπεδα ακεραιότητας στα Windows}{Επίπεδο: Εισαγωγικό–Μεσαίο επίπεδο · Συστήματα}{Χρόνος μελέτης: 7–9 ώρες}

Το λειτουργικό σύστημα είναι το πρώτο απτό όριο ασφάλειας που συναντά ο φοιτητής. Κάθε αρχείο, κάθε διεργασία και κάθε δικτυακή υποδοχή (socket) ελέγχεται τελικά από τον πυρήνα, ο οποίος αποφασίζει ποιος μπορεί να κάνει τι. Αν ο έλεγχος πρόσβασης του λειτουργικού συστήματος είναι λανθασμένα διαμορφωμένος, οι άμυνες υψηλότερου επιπέδου, όπως τα τείχη προστασίας ή το λογισμικό προστασίας από ιούς, συχνά παρακάμπτονται με ελάχιστη προσπάθεια.

Το κεφάλαιο συγκρίνει τις δύο κυρίαρχες οικογένειες ελέγχου πρόσβασης. Αρχικά εξηγεί το μοντέλο POSIX που χρησιμοποιούν το Linux, το macOS και τα υπόλοιπα συστήματα τύπου Unix, συμπεριλαμβανομένων των ειδικών δικαιωμάτων που αποτελούν συχνή αιτία κλιμάκωσης προνομίων. Στη συνέχεια προχωρά πέρα από τα απλά δικαιώματα, στις λίστες ελέγχου πρόσβασης (ACL), στις δυνατότητες (capabilities) του Linux και στα πλαίσια υποχρεωτικού ελέγχου πρόσβασης. Το δεύτερο μισό αφιερώνεται στα Windows: αναγνωριστικά ασφάλειας, διακριτικά πρόσβασης (access tokens), Έλεγχος Λογαριασμού Χρήστη (UAC), το μητρώο (registry) και τα υποχρεωτικά επίπεδα ακεραιότητας.

\begin{outcomesbox}{Μαθησιακά αποτελέσματα}
Μετά τη μελέτη του κεφαλαίου θα πρέπει να μπορείτε:
\begin{itemize}
  \item Να διαβάζετε, να ερμηνεύετε και να ορίζετε δικαιώματα POSIX σε συμβολική και οκταδική μορφή.
  \item Να εξηγείτε τον σκοπό και τους κινδύνους των ειδικών δικαιωμάτων SetUID, SetGID και sticky bit.
  \item Να περιγράφετε πώς οι ACL, οι δυνατότητες (capabilities) και ο υποχρεωτικός έλεγχος πρόσβασης (SELinux, AppArmor) εξειδικεύουν το βασικό μοντέλο δικαιωμάτων.
  \item Να εξηγείτε τον ρόλο των SID, των διακριτικών πρόσβασης, του UAC και των επιπέδων ακεραιότητας στα Windows.
  \item Να εντοπίζετε συνήθεις εσφαλμένες ρυθμίσεις λειτουργικών συστημάτων που επιτρέπουν κλιμάκωση προνομίων ή πλευρική κίνηση.
\end{itemize}
\end{outcomesbox}

\section{Διακριτικός έλεγχος πρόσβασης POSIX: μοντέλο και λειτουργία}\label{sec:2.1}
Τα συστήματα POSIX εφαρμόζουν \textbf{διακριτικό έλεγχο πρόσβασης} (Discretionary Access Control, DAC): ο κάτοχος ενός αρχείου αποφασίζει ποιοι άλλοι μπορούν να έχουν πρόσβαση σε αυτό. Κάθε αρχείο έχει έναν κάτοχο (αναγνωριστικό χρήστη), μια ομάδα (αναγνωριστικό ομάδας) και τρία σύνολα δικαιωμάτων —για τον \emph{κάτοχο}, την \emph{ομάδα} και τους \emph{λοιπούς}. Κάθε σύνολο περιλαμβάνει τρία δικαιώματα: \textbf{ανάγνωση (r)}, \textbf{εγγραφή (w)} και \textbf{εκτέλεση (x)}. Στους καταλόγους η σημασία διαφοροποιείται ελαφρώς: η ανάγνωση επιτρέπει την εμφάνιση των ονομάτων, η εγγραφή τη δημιουργία ή διαγραφή εγγραφών και η εκτέλεση την είσοδο στον κατάλογο και την πρόσβαση στα αρχεία που περιέχει.

Τα δικαιώματα γράφονται συνήθως σε οκταδική μορφή, όπου η ανάγνωση αντιστοιχεί στο 4, η εγγραφή στο 2 και η εκτέλεση στο 1. Επομένως, η τιμή \texttt{640} σημαίνει \emph{rw-} για τον κάτοχο, \emph{r--} για την ομάδα και \emph{---} για όλους τους υπόλοιπους —μια τυπική ρύθμιση για αρχείο ρυθμίσεων που περιέχει μυστικά. Όταν μια διεργασία ζητά πρόσβαση, ο πυρήνας ελέγχει τις κατηγορίες με σειρά (κάτοχος, ομάδα, λοιποί) και εφαρμόζει μόνο την \emph{πρώτη} κατηγορία που ταιριάζει. Κατά συνέπεια, ένας κάτοχος που έχει αφαιρέσει από τον εαυτό του το δικαίωμα ανάγνωσης δεν αποκτά πρόσβαση, ακόμη κι αν οι «λοιποί» μπορούν να διαβάσουν το αρχείο.

\begin{examplebox}{Παράδειγμα εφαρμογής}
Η εντολή \texttt{chmod 750 /srv/reports} παρέχει στον κάτοχο πλήρη έλεγχο (7 = 4+2+1), επιτρέπει στα μέλη της ομάδας να βλέπουν και να εισέρχονται στον κατάλογο (5 = 4+1) και απαγορεύει κάθε πρόσβαση στους υπόλοιπους χρήστες (0). Ένας διακομιστής ιστού που εκτελείται με άσχετο λογαριασμό δεν θα μπορεί, επομένως, να διαβάσει τις αναφορές, ακόμη κι αν μια ευπάθεια του επιτρέπει να μαντέψει τη διαδρομή τους.
\end{examplebox}

\section{Ειδικά δικαιώματα: SetUID, SetGID και sticky bit}\label{sec:2.2}
Τρία πρόσθετα bit τροποποιούν τη συμπεριφορά των βασικών δικαιωμάτων. Το \textbf{SetUID} (οκταδικό 4000) προκαλεί την εκτέλεση ενός προγράμματος με τα προνόμια του \emph{κατόχου} του και όχι του χρήστη που το εκκίνησε. Με αυτόν τον τρόπο ένας απλός χρήστης μπορεί να αλλάξει τον κωδικό του με την εντολή \texttt{passwd}, ένα πρόγραμμα που ανήκει στον root και πρέπει να γράψει στο προστατευμένο αρχείο \texttt{/etc/shadow}. Το \textbf{SetGID} (2000) λειτουργεί ανάλογα για την ομάδα· όταν εφαρμόζεται σε κατάλογο, κάνει επιπλέον τα νέα αρχεία να κληρονομούν την ομάδα του καταλόγου, κάτι χρήσιμο για κοινόχρηστους φακέλους έργων. Το \textbf{sticky bit} (1000) σε κατάλογο με δικαίωμα εγγραφής για όλους, όπως ο \texttt{/tmp}, επιτρέπει σε κάθε χρήστη να διαγράφει μόνο τα δικά του αρχεία.

Τα προγράμματα SetUID που ανήκουν στον root είναι από τους πιο ελκυστικούς στόχους για τους επιτιθέμενους, επειδή οποιοδήποτε σφάλμα τους —υπερχείλιση ενδιάμεσης μνήμης, μη ασφαλής κλήση άλλου προγράμματος, διαδρομή βιβλιοθήκης με δικαίωμα εγγραφής— οδηγεί αμέσως σε πλήρη διαχειριστικό έλεγχο. Γι' αυτό οι διαχειριστές συστημάτων οφείλουν να τα καταγράφουν περιοδικά (για παράδειγμα με την εντολή \texttt{find / -perm -4000 -type f}), να αφαιρούν το bit όπου δεν είναι απολύτως αναγκαίο και να προτιμούν λεπτομερέστερους μηχανισμούς, όπως οι δυνατότητες (capabilities).

\section{Πέρα από τα βασικά δικαιώματα: ACL, capabilities και υποχρεωτικός έλεγχος πρόσβασης}\label{sec:2.3}
Το μοντέλο κάτοχος–ομάδα–λοιποί είναι απλό, αλλά δεν μπορεί να εκφράσει κανόνες όπως \emph{η Αλίκη μπορεί να διαβάσει αυτό το αρχείο και ο Βασίλης να το τροποποιήσει, ενώ κανείς άλλος δεν μπορεί να κάνει τίποτα από τα δύο}. Οι \textbf{λίστες ελέγχου πρόσβασης POSIX (ACL)} επιλύουν το πρόβλημα, προσθέτοντας εγγραφές για συγκεκριμένους χρήστες και ομάδες, οι οποίες διαχειρίζονται με την εντολή \texttt{setfacl} και εμφανίζονται με την \texttt{getfacl}. Μια εγγραφή \emph{μάσκας} περιορίζει τα μέγιστα δικαιώματα που μπορεί να λάβει οποιαδήποτε ονομαστική εγγραφή, επιτρέποντας στους διαχειριστές να περιορίζουν γρήγορα την πρόσβαση.

Οι \textbf{δυνατότητες (capabilities) του Linux} διασπούν το παντοδύναμο προνόμιο του root σε περίπου σαράντα ανεξάρτητες μονάδες. Ένας διακομιστής ιστού που χρειάζεται μόνο να δεσμεύσει τη θύρα 80 μπορεί να λάβει τη δυνατότητα \texttt{CAP\_NET\_BIND\_SERVICE} αντί να εκτελείται ως root, ώστε η παραβίασή του να μην παρέχει έλεγχο ολόκληρου του συστήματος. Τέλος, τα πλαίσια \textbf{υποχρεωτικού ελέγχου πρόσβασης} (Mandatory Access Control, MAC), όπως τα \textbf{SELinux} και \textbf{AppArmor}, επιβάλλουν μια πολιτική σε επίπεδο συστήματος την οποία δεν μπορεί να παρακάμψει ούτε ο κάτοχος του αρχείου. Υπό καθεστώς MAC, μια παραβιασμένη διεργασία παραμένει περιορισμένη στους πόρους που της επιτρέπει ρητά η πολιτική —εφαρμογή της αρχής των ασφαλών προεπιλογών που εισήχθη στο Κεφάλαιο 1.

\section{Αρχιτεκτονική ασφάλειας των Windows: ταυτότητες, token, UAC και μητρώο}\label{sec:2.4}
Τα Windows αναγνωρίζουν κάθε χρήστη, ομάδα και υπολογιστή μέσω ενός \textbf{αναγνωριστικού ασφάλειας} (Security Identifier, SID), μιας μοναδικής τιμής που παραμένει σταθερή ακόμη κι αν ο λογαριασμός μετονομαστεί. Όταν ένας χρήστης συνδέεται, η Τοπική Αρχή Ασφάλειας (LSA) δημιουργεί ένα \textbf{διακριτικό πρόσβασης} (access token), το οποίο περιέχει το SID του χρήστη, τα SID των ομάδων του και τα προνόμιά του. Κάθε διεργασία που εκκινεί ο χρήστης κληρονομεί αντίγραφο αυτού του διακριτικού. Αντικείμενα όπως αρχεία, κλειδιά μητρώου και υπηρεσίες φέρουν έναν \textbf{περιγραφέα ασφάλειας} (security descriptor), του οποίου η \textbf{λίστα διακριτικού ελέγχου πρόσβασης} (DACL) περιέχει εγγραφές που επιτρέπουν ή απαγορεύουν την πρόσβαση. Όταν μια διεργασία ζητά πρόσβαση, η Οθόνη Αναφοράς Ασφάλειας (Security Reference Monitor) συγκρίνει το διακριτικό με τη DACL· οι ρητές απαγορεύσεις αξιολογούνται πριν από τις άδειες.

Ο \textbf{Έλεγχος Λογαριασμού Χρήστη} (User Account Control, UAC) αντιμετωπίζει ένα ιστορικό πρόβλημα: για χρόνια, οι περισσότεροι χρήστες των Windows εργάζονταν μόνιμα ως διαχειριστές. Με το UAC, ένας διαχειριστής λαμβάνει κατά τη σύνδεση δύο διακριτικά —ένα \emph{φιλτραρισμένο} διακριτικό απλού χρήστη για την καθημερινή εργασία και ένα \emph{πλήρες} διακριτικό που ενεργοποιείται μόνο μετά από ρητή συγκατάθεση. Το \textbf{μητρώο} (registry), η ιεραρχική βάση δεδομένων που αποθηκεύει τις ρυθμίσεις του συστήματος και των εφαρμογών, προστατεύεται από τους ίδιους περιγραφείς ασφάλειας. Κλειδιά μητρώου όπως το \texttt{HKLM\textbackslash{}...\textbackslash{}Run} αποτελούν αγαπημένο σημείο εγκατάστασης μόνιμης παρουσίας (persistence) για το κακόβουλο λογισμικό· γι' αυτό τα δικαιώματα και οι μεταβολές τους πρέπει να παρακολουθούνται.

\section{Τα Windows σε βάθος: επίπεδα ακεραιότητας, προνόμια και επιφάνεια πλευρικής κίνησης}\label{sec:2.5}
Από τα Windows Vista και έπειτα, κάθε διεργασία και κάθε αντικείμενο φέρει επιπλέον ένα \textbf{υποχρεωτικό επίπεδο ακεραιότητας}: \emph{Χαμηλό}, \emph{Μεσαίο}, \emph{Υψηλό} ή \emph{Συστήματος}. Μια διεργασία δεν μπορεί να γράψει σε αντικείμενο υψηλότερου επιπέδου ακεραιότητας, ανεξάρτητα από το τι ορίζει η DACL. Πρόκειται για άμεση εφαρμογή του μοντέλου Biba (Κεφάλαιο 1). Οι φυλλομετρητές ιστού αξιοποιούν αυτόν τον μηχανισμό, εκτελώντας τις διεργασίες απόδοσης περιεχομένου σε Χαμηλό επίπεδο ακεραιότητας, ώστε κακόβουλο περιεχόμενο ιστού να μην μπορεί να τροποποιήσει τα αρχεία του χρήστη ακόμη κι αν επιτύχει εκτέλεση κώδικα.

Ορισμένα \textbf{προνόμια} ενός διακριτικού είναι τόσο ισχυρά, ώστε ισοδυναμούν ουσιαστικά με πλήρη έλεγχο. Το \texttt{SeDebugPrivilege} επιτρέπει την ανάγνωση της μνήμης οποιασδήποτε διεργασίας —συμπεριλαμβανομένης της LSASS, όπου διατηρούνται διαπιστευτήρια— ενώ τα \texttt{SeImpersonatePrivilege} και \texttt{SeBackupPrivilege} έχουν επανειλημμένα χρησιμοποιηθεί για κλιμάκωση προνομίων. Σε ένα περιβάλλον τομέα (domain), αυτές οι τοπικές αδυναμίες μετατρέπονται σε δικτυακό πρόβλημα: τα αποθηκευμένα διαπιστευτήρια, οι κοινοί κωδικοί τοπικού διαχειριστή και τα πρωτόκολλα απομακρυσμένης διαχείρισης (SMB, WMI, WinRM, RDP) επιτρέπουν σε έναν επιτιθέμενο να κινηθεί πλευρικά από έναν παραβιασμένο σταθμό εργασίας σε πολλούς άλλους. Τα αντίμετρα περιλαμβάνουν μοναδικούς κωδικούς τοπικού διαχειριστή (Windows LAPS), το Credential Guard, την κλιμακωτή (tiered) διαχείριση και τον περιορισμό των λογαριασμών που επιτρέπεται να συνδέονται σε κάθε σύστημα.

\begin{table}[htbp]
  \centering\small
  \caption{Σύγκριση ελέγχου πρόσβασης σε POSIX και Windows}
  \begin{tabularx}{\linewidth}{>{\raggedright\arraybackslash}X >{\raggedright\arraybackslash}X >{\raggedright\arraybackslash}X}
    \toprule
    \textbf{Χαρακτηριστικό} & \textbf{POSIX (Linux)} & \textbf{Windows} \\
    \midrule
    Ταυτότητα & UID / GID & SID \\
    Πολιτική ανά αντικείμενο & Βασικά δικαιώματα + προαιρετική ACL & Περιγραφέας ασφάλειας με DACL \\
    Ανύψωση προνομίων & SetUID, sudo, capabilities & Συγκατάθεση UAC, προνόμια στο διακριτικό \\
    Υποχρεωτικός έλεγχος & SELinux, AppArmor & Επίπεδα ακεραιότητας, AppContainer \\
    Συνήθης οδός κλιμάκωσης & Ευάλωτο εκτελέσιμο SetUID & Κατάχρηση προνομίου διακριτικού, αδύναμη ACL υπηρεσίας \\
    \bottomrule
  \end{tabularx}
\end{table}
\section*{Βασικοί όροι}
\begin{description}[style=nextline,leftmargin=1.2em]
  \item[DAC] Διακριτικός έλεγχος πρόσβασης: ο κάτοχος ενός πόρου αποφασίζει ποιος έχει πρόσβαση σε αυτόν.
  \item[SetUID] Ειδικό δικαίωμα που κάνει ένα πρόγραμμα να εκτελείται με τα προνόμια του κατόχου του.
  \item[Δυνατότητα (capability)] Αυτοτελής μονάδα του προνομίου root που μπορεί να παραχωρηθεί χωριστά.
  \item[MAC] Υποχρεωτικός έλεγχος πρόσβασης: πολιτική σε επίπεδο συστήματος που δεν μπορεί να παρακάμψει ο κάτοχος (SELinux, AppArmor).
  \item[Διακριτικό πρόσβασης] Δομή των Windows που περιέχει την ταυτότητα, τις ομάδες και τα προνόμια μιας διεργασίας.
  \item[Επίπεδο ακεραιότητας] Ετικέτα των Windows (Χαμηλό–Συστήματος) που εμποδίζει διεργασίες χαμηλότερου επιπέδου να γράφουν σε αντικείμενα υψηλότερου.
\end{description}

\section*{Σύνοψη κεφαλαίου}
\begin{itemize}
  \item Τα συστήματα POSIX ελέγχουν την πρόσβαση μέσω δικαιωμάτων για κάτοχο, ομάδα και λοιπούς· αποφασίζει η πρώτη κατηγορία που ταιριάζει.
  \item Τα SetUID και SetGID είναι αναγκαία σε λίγα σημεία, αλλά αποτελούν βασικό στόχο κλιμάκωσης προνομίων και πρέπει να καταγράφονται.
  \item Οι ACL προσθέτουν λεπτομερείς διακριτικούς κανόνες, οι δυνατότητες διασπούν το προνόμιο του root και ο MAC περιορίζει ακόμη και παραβιασμένες διεργασίες.
  \item Τα Windows συνδυάζουν SID, διακριτικά πρόσβασης και DACL· το UAC διαχωρίζει την καθημερινή εργασία από τις διαχειριστικές ενέργειες.
  \item Τα επίπεδα ακεραιότητας και η προσεκτική διαχείριση προνομίων περιορίζουν τόσο την τοπική κλιμάκωση όσο και την πλευρική κίνηση σε έναν τομέα.
\end{itemize}

\section*{Ερωτήσεις ανασκόπησης}
\begin{enumerate}
  \item Ένα αρχείο έχει δικαιώματα \texttt{604} και κάτοχο την alice. Μπορεί η alice να το διαβάσει; Μπορεί ο bob; Εξηγήστε τη σειρά αξιολόγησης που οδηγεί στην απάντησή σας.
  \item Γιατί ένα ευάλωτο εκτελέσιμο SetUID που ανήκει στον root είναι πιο επικίνδυνο από την ίδια ευπάθεια σε ένα συνηθισμένο πρόγραμμα;
  \item Συγκρίνετε τις δυνατότητες του Linux με τα προνόμια διακριτικού των Windows. Ποια αρχή ασφάλειας υλοποιούν και οι δύο μηχανισμοί;
  \item Εξηγήστε πώς οι κοινοί κωδικοί τοπικού διαχειριστή διευκολύνουν την πλευρική κίνηση και περιγράψτε δύο αντίμετρα.
\end{enumerate}
